\section{Discussion and conclusion}\label{sec:discussion}

\subsection{What improves a usable design certificate}

The distinction between a statistical information model and an optimization
relaxation determines what can be certified. The public-source enumeration
shows that full-inverse gating can change a selected experiment even in a
small exactly enumerable model. It also shows why the scope of a correction
must be precise: D-optimal and conventional A-optimal choices agree across all
eleven tested budgets, while one trace-information choice changes. The
unconditional information ordering in Proposition~\ref{prop:gating} explains
how a gated global upper bound can still be reused after the incumbent is
evaluated under selected covariance. The condition-dependent comparison
separately bounds the possible efficiency loss.
Agreement of choices alone does not establish equality of information values.

For temporal designs, finite histories are useful when their conditional
structure supports a tractable full-family price and their uniform error is
small enough for the desired certificate. The positive all-split comparisons
are stronger evidence of relaxation strength than a timed solver contest.
Their common fractional witnesses prove that better continuous optimization
or scalar split selection cannot eliminate the measured separations. The
kinetic all-diagonal witness extends that conclusion to a larger family and
its specified affine upper supports. It does not limit branching or other
valid cuts. Conversely, the complete-packet examples show that the tested
dense bound can be tighter and that all methods can select the same design.
There is no demonstrated universal dominance.

The results suggest allocating work according to the part of the bound that
remains loose. A large continuous support residual calls for a better reference
or further numerical optimization. A large locality correction calls for a
stronger model-specific bound, a longer history, or a different representation.
A remaining mixture integrality gap requires restrictions that distinguish
one feasible experiment from a fractional mixture. More accurate arithmetic
can reduce enclosure loss, as the spacing experiment illustrates, but cannot
remove either a statistical truncation bound or a mixture gap. This separation
also permits inexpensive heuristics to supply incumbents while independently
chosen witnesses certify them.

Latent separators expose a second tradeoff. For nested anchors and the same
complete feasible family, removing anchors improves the mathematical hull
bound. The exact eight-time study verifies strict improvements rather than
comparing unrelated partitions or unfinished solver values. At the same time,
fewer anchors require larger conditional-block patterns, while more anchors
increase nuisance dimension. The larger fixed-physical experiments retain
these costs: the best reported 192-point separator certificate still misses
the .01 log-gap target and exceeds the nominal 30-second allowance once
certification is included. The hierarchy establishes an ordering of bounds,
not a uniformly efficient choice of representation.

\subsection{Scope of the statistical and complexity conclusions}

Fixed covariance and nominal sensitivities make the information objective
precise, but limit its interpretation. The kinetic examples are local designs
for stipulated error models. Independent balance-equation and high-precision
checks support their sensitivity implementation; the exact certificates apply
to the stored rational arrays. Neither step validates a physical noise law.
The exact rate-swap symmetry also persists despite full-rank local information.
A positive definite regularizer supplies the declared algebraic prior term,
rather than resolving this global ambiguity.

Finite-scenario standardization answers a different question from raw maximin
information. It measures each scenario against its own attainable optimum,
and its certificate must therefore include uncertainty in those optima. The
shared support bound prices one design over all scenarios and improves on
separate scenario caps in the reported tests. Its guarantee is restricted to
the three listed regimes. Covering a continuous parameter region would require
additional bounds between scenarios; the present experiment does not assert
such coverage.

The approximation theorems identify tractable promised classes rather than
predicting the runtime of the numerical prototypes. Weighted trace permits
growing packet and parameter dimensions because its local support is linear.
Nonlinear criteria use a fixed-dimensional relative PSD cover with large
polynomial exponents and explicit feasibility representations. The correlated
temporal schemes require complete packets and the stated fixed covariance
promises. The abstract additive cover instead requires fixed information
dimension and the specified explicit rational feasibility representation.
For A- and E-design with additive rank-one information under partition constraints,
\citet{BansalXu2026} show that, when information dimension grows, approximation
within any polynomial factor in that dimension is impossible in polynomial
time unless $\mathsf{P}=\mathsf{NP}$, even when every feasible information matrix
is invertible. This is consistent with the fixed-dimensional cover here.
A finer grid for one physical process makes the per-step
correlation approach one; the measured calendar-price failures are consistent
with this loss of uniform promises. The conservative sufficient windows do
not prove a lower bound on every possible algorithm. Likewise, the result for
represented matroids concerns additive information atoms and provides no automatic
way to impose arbitrary matroid constraints on correlated histories.

\subsection{Conclusion}

The central result is a route from a specified correlated experiment to a
checkable global design bound. Explicit local conditional errors support
finite-state certificates and approximation guarantees; normalization on exact
ranges extends the latter to several matrix criteria; and Schur
elimination supplies a nested alternative when calendar histories are costly.
The computations establish strict relaxation-family separations and useful
finite-scenario guarantees, while retaining cases with weaker bounds or
unfavorable cost. Their common standard is a feasible experiment and a verified
upper witness for the same selected-covariance objective. This makes the
guarantees inspectable independently of the numerical method that proposed
the design.

\subsection*{Data and code availability}

The accompanying standalone computational supplement contains the model arrays,
selected schedules, rational witnesses, exact checkers, numerical producers,
dependency lock, and file-hash manifests used in this paper. Its README maps
the reported studies and tables to the saved artifacts and gives separate
commands for certificate validation and numerical reproduction. Full validation
replays 46 principal saved certificates, independently checks every public-source
schedule, checks kinetic sensitivities, and verifies the fresh exhaustive
studies. After dependencies are installed, validation requires no network,
commercial solver, original repository, or literature collection. Optional
historical mixed-integer proposal reproduction uses Gurobi and its applicable
license. The public sensitivity CSV retains its immutable source reference
\citep{MeasurementOpt2025} and hash; third-party literature originals and
uninspected stored author selections are not included.
